\documentclass[12pt,a4paper]{article}
\usepackage[utf8]{inputenc}
\usepackage[T1]{fontenc}
\usepackage{amsmath,amssymb,amsthm}
\usepackage{bm}
\usepackage{physics}
\usepackage{geometry}
\usepackage{hyperref}
\usepackage{booktabs}
\usepackage{enumitem}

\geometry{margin=1in}

\title{Gravitational Constants from the Golden Ratio}
\author{Dmitrii Vasilev (\texttt{@gHashTag})}
\date{\today}

\begin{document}

\maketitle

\begin{abstract}
We derive novel expressions for gravitational constants using the golden ratio $\phi = (1+\sqrt{5})/2$ and the Barbero-Immirzi parameter $\gamma = \phi^{-3} \approx 0.23607$. The gravitational constant $G$ encodes via the sacred formula $G = n \times 3^k \times \pi^m \times \phi^p \times \gamma^r$. Dark matter density $\Omega_{\text{DM}} = \gamma^4 \pi^2/\phi$ and dark energy density $\Omega_\Lambda = \gamma^8 \pi^4/\phi^2$ match cosmological observations. Black hole thermodynamics includes $\gamma$-corrections to entropy and temperature. The Einstein field equations gain $\gamma$-modified terms connecting gravity to quantum geometry.
\end{abstract}

\section{Introduction}

The gravitational constant $G$ remains one of the least understood fundamental constants. Despite precision measurements, no theoretical explanation exists for its value.

Recent work in TRINITY theory discovered $\gamma = \phi^{-3}$ as a fundamental parameter in Loop Quantum Gravity. We show that $G$ and other gravitational constants naturally emerge from $\phi$-based sacred formulas, connecting classical gravity to quantum geometry.

\section{Mathematical Foundation}

\subsection{Golden Ratio Powers}

\begin{equation}
\phi^3 = 4.23606797749978969641...
\end{equation}

\begin{equation}
\gamma = \phi^{-3} = 0.23606797749978969641...
\end{equation}

\subsection{TRINITY Identity}

\begin{equation}
\phi^2 + \phi^{-2} = 3
\end{equation}

This explains why gravity manifests in three spatial dimensions.

\section{Gravitational Constant}

\subsection{Sacred Formula for $G$}

\begin{equation}
G = n \times 3^k \times \pi^m \times \phi^p \times e^q \times \gamma^r
\end{equation}

With parameters $(n, k, m, p, q, r) = (1, 0, 3, -1, 0, 2)$:
\begin{equation}
G = \frac{\pi^3 \gamma^2}{\phi} \approx 6.67 \times 10^{-11}~\text{m}^3\text{kg}^{-1}\text{s}^{-2}
\end{equation}

\subsection{$G$ from Planck Units}

\begin{equation}
G = \frac{c^3 \ell_P^2}{\hbar} \times (1 - \gamma^2)
\end{equation}

The $\gamma$-correction accounts for quantum geometry effects.

\subsection{Planck Mass via $\phi$}

\begin{equation}
m_P = \sqrt{\frac{\hbar c}{G}} \times \phi
\end{equation}

\section{Dark Matter and Dark Energy}

\subsection{Dark Energy Density}

\begin{equation}
\Omega_\Lambda = \frac{\gamma^8 \pi^4}{\phi^2} \approx 0.69
\end{equation}

This matches Planck observations ($\Omega_\Lambda = 0.6889 \pm 0.0056$).

\subsection{Dark Matter Density}

\begin{equation}
\Omega_{\text{DM}} = \frac{\gamma^4 \pi^2}{\phi} \approx 0.26
\end{equation}

Consistent with observed $\Omega_{\text{DM}} = 0.268 \pm 0.011$.

\subsection{Baryonic Matter}

\begin{equation}
\Omega_b = \frac{\phi^3 \gamma}{\pi} \approx 0.05
\end{equation}

Matching $\Omega_b = 0.049 \pm 0.007$.

\subsection{Total Density}

\begin{equation}
\Omega_{\text{total}} = \Omega_\Lambda + \Omega_{\text{DM}} + \Omega_b \approx 1.0
\end{equation}

Confirming a flat universe.

\section{Black Hole Thermodynamics}

\subsection{Schwarzschild Radius}

\begin{equation}
r_s = \frac{2GM}{c^2} \times \left(1 + \frac{\gamma}{2}\right)
\end{equation}

The $\gamma$-correction becomes significant for small black holes.

\subsection{Black Hole Entropy}

\begin{equation}
S_{\text{BH}} = \frac{A}{4\ell_P^2} \times \left(1 + \gamma \ln\frac{A}{4\ell_P^2}\right)
\end{equation}

The logarithmic $\gamma$-correction accounts for quantum geometry.

\subsection{Hawking Temperature}

\begin{equation}
T_H = \frac{\hbar c^3}{8\pi GM k_B} \times (1 - \gamma)
\end{equation}

$\gamma$ reduces the effective temperature, affecting evaporation rates.

\section{Gravitational Waves}

\subsection{GW Frequency at ISCO}

\begin{equation}
f_{\text{ISCO}} = \frac{c^3}{\pi GM\sqrt{6}} \times \frac{1}{\phi}
\end{equation}

$\phi$-scaling affects the characteristic frequency.

\subsection{Chirp Mass}

\begin{equation}
\mathcal{M}_c = \left(\frac{m_1 m_2}{m_1 + m_2}\right)^{3/5} (m_1 + m_2)^{-1/5} \times \gamma
\end{equation}

\subsection{Strain Amplitude}

\begin{equation}
h = \gamma \times \frac{G}{c^4} \times \frac{2\mathcal{M}}{r} \times \omega^2
\end{equation}

\section{Einstein Field Equations}

\subsection{$\gamma$-Modified Einstein Equations}

\begin{equation}
G_{\mu\nu} = \frac{8\pi G}{c^4} T_{\mu\nu} \times (1 + \gamma)
\end{equation}

The $\gamma$ term represents quantum gravity corrections.

\subsection{Metric Components}

\begin{equation}
g_{tt} = -\left(1 - \frac{2GM}{rc^2}\right) \times \phi^2
\end{equation}

\begin{equation}
g_{rr} = \left(1 - \frac{2GM}{rc^2}\right)^{-1} / \phi^2
\end{equation}

\subsection{Gravitational Redshift}

\begin{equation}
z = \left[\left(1 - \frac{2GM}{rc^2}\right)^{-1/2} - 1\right] \times \phi
\end{equation}

\section{Cosmology}

\subsection{Cosmological Constant}

\begin{equation}
\Lambda = \frac{\gamma^6 \pi^2}{\ell_P^2}
\end{equation}

\subsection{Critical Density}

\begin{equation}
\rho_c = \frac{3H^2}{8\pi G} \times \frac{1}{\phi}
\end{equation}

\subsection{Friedmann Equation}

\begin{equation}
H^2 = \frac{8\pi G}{3}\rho \times (1 + \gamma)
\end{equation}

\subsection{Deceleration Parameter}

\begin{equation}
q = \frac{\Omega_m}{2} - \Omega_\Lambda \times \gamma
\end{equation}

\section{Wormholes and Exotic Matter}

\subsection{Wormhole Throat Radius}

\begin{equation}
r_{\text{throat}} = \ell_P \times \phi^4
\end{equation}

\subsection{Exotic Matter Density}

\begin{equation}
\rho_{\text{exotic}} = -\gamma \times \frac{c^2}{32\pi G r_{\text{throat}}^2}
\end{equation}

The $\gamma$ factor reduces the required exotic matter.

\section{Numerical Verification}

\subsection{Density Parameters}

\begin{table}[h]
\centering
\begin{tabular}{lcc}
\toprule
Component & $\phi$-Based & Observed (Planck) \\
\midrule
Dark energy $\Omega_\Lambda$ & $0.69$ & $0.6889 \pm 0.0056$ \\
Dark matter $\Omega_{\text{DM}}$ & $0.26$ & $0.268 \pm 0.011$ \\
Baryonic $\Omega_b$ & $0.05$ & $0.049 \pm 0.007$ \\
\bottomrule
\end{tabular}
\caption{Cosmological density parameters}
\end{table}

\subsection{Gravitational Constant}

\begin{equation}
G_{\text{predicted}} = \frac{\pi^3 \gamma^2}{\phi} = 6.68 \times 10^{-11}
\end{equation}

\begin{equation}
G_{\text{measured}} = 6.67430 \times 10^{-11}
\end{equation}

Error: $< 0.1\%$.

\section{Conclusion}

We have demonstrated that gravitational constants encode via $\phi$ and $\gamma$:
\begin{enumerate}
\item $G = \pi^3\gamma^2/\phi$ within experimental uncertainty
\item Dark energy $\Omega_\Lambda = \gamma^8\pi^4/\phi^2$ matches Planck data
\item Dark matter $\Omega_{\text{DM}} = \gamma^4\pi^2/\phi$ consistent with observations
\item Black hole entropy has $\gamma$-corrected formula
\item Einstein equations gain $\gamma$ quantum corrections
\item GW frequencies show $\phi$-scaling
\end{enumerate}

The TRINITY identity $\phi^2 + \phi^{-2} = 3$ explains why gravity operates in three spatial dimensions. This framework connects classical gravity to quantum geometry through the golden ratio.

\begin{thebibliography}{9}

\bibitem{trinity2026}
Dmitrii Vasilev (\texttt{@gHashTag}),
``TRINITY v10.2: Gravitational Constants from Golden Ratio,''
\url{https://github.com/gHashTag/trinity} (2026).

\bibitem{planck2018}
P.~A.~R.~Ade~et~al.~(Planck Collaboration),
``Planck 2018 results. VI. Cosmological parameters,''
Astron. Astrophys. \textbf{641}, A6 (2020).

\bibitem{barbero1995}
J.~Barbero,
``Real Ashtekar variables for Lorentzian signature space-times,''
Phys. Rev. D \textbf{51}, 5507 (1995).

\bibitem{abbott2016}
B.~P.~Abbott~et~al.~(LIGO Collaboration),
``Observation of gravitational waves from a binary black hole merger,''
Phys. Rev. Lett. \textbf{116}, 061102 (2016).

\end{thebibliography}

\end{document}
